\documentclass[a4paper,12pt]{article}

% ===========================================================================
%  Лабораторная работа 3.4.5 «Петля гистерезиса (динамический метод)»
%
%  Сборка:  cd 3.4.5/report && latexmk -pdf main.tex
%  Таблицы и графики читаются из ../data/processed/*.csv, итоговые числа --
%  из ../data/processed/results.tex (макрос \res{...}). Всё это готовит
%  ../scripts/process_data.py. До проведения измерений на месте чисел
%  стоят метки ???.
% ===========================================================================

% ============================================================
% Общий преамбул для отчётов по лабораторным работам (labs_lpr)
% Подключается из report/main.tex командой:
%   % ============================================================
% Общий преамбул для отчётов по лабораторным работам (labs_lpr)
% Подключается из report/main.tex командой:
%   % ============================================================
% Общий преамбул для отчётов по лабораторным работам (labs_lpr)
% Подключается из report/main.tex командой:
%   \input{../../common/latex_template/preamble.tex}
% ============================================================

% Поддержка русского языка и кодировки
\usepackage[T2A]{fontenc}
\usepackage{lmodern}
\usepackage[utf8]{inputenc}
\usepackage[russian]{babel}

% Математические пакеты
\usepackage{amsmath, amssymb}

% Настройка полей
\usepackage{geometry}
\geometry{top=2cm, bottom=2cm, left=2.5cm, right=1.5cm}

% Колонтитулы
\usepackage{fancyhdr}
\usepackage{graphicx}

% Современные таблицы
\usepackage{tabularray}

% Работа с csv файлами
\usepackage{csvsimple-l3}
\usepackage{tikz}

% Графики
\usepackage{pgfplots}
\pgfplotsset{compat=1.18}

% Ссылки
% \usepackage{hyperref} % гипперссылки, нужно оформление вместо красного квадрата
\usepackage[russian]{cleveref} % умные ссылки, нужно переоформление для формул, рисунков и тд (вместо ф-л что-то нужно придумать)

% ============================================================

% Поддержка русского языка и кодировки
\usepackage[T2A]{fontenc}
\usepackage{lmodern}
\usepackage[utf8]{inputenc}
\usepackage[russian]{babel}

% Математические пакеты
\usepackage{amsmath, amssymb}

% Настройка полей
\usepackage{geometry}
\geometry{top=2cm, bottom=2cm, left=2.5cm, right=1.5cm}

% Колонтитулы
\usepackage{fancyhdr}
\usepackage{graphicx}

% Современные таблицы
\usepackage{tabularray}

% Работа с csv файлами
\usepackage{csvsimple-l3}
\usepackage{tikz}

% Графики
\usepackage{pgfplots}
\pgfplotsset{compat=1.18}

% Ссылки
% \usepackage{hyperref} % гипперссылки, нужно оформление вместо красного квадрата
\usepackage[russian]{cleveref} % умные ссылки, нужно переоформление для формул, рисунков и тд (вместо ф-л что-то нужно придумать)

% ============================================================

% Поддержка русского языка и кодировки
\usepackage[T2A]{fontenc}
\usepackage{lmodern}
\usepackage[utf8]{inputenc}
\usepackage[russian]{babel}

% Математические пакеты
\usepackage{amsmath, amssymb}

% Настройка полей
\usepackage{geometry}
\geometry{top=2cm, bottom=2cm, left=2.5cm, right=1.5cm}

% Колонтитулы
\usepackage{fancyhdr}
\usepackage{graphicx}

% Современные таблицы
\usepackage{tabularray}

% Работа с csv файлами
\usepackage{csvsimple-l3}
\usepackage{tikz}

% Графики
\usepackage{pgfplots}
\pgfplotsset{compat=1.18}

% Ссылки
% \usepackage{hyperref} % гипперссылки, нужно оформление вместо красного квадрата
\usepackage[russian]{cleveref} % умные ссылки, нужно переоформление для формул, рисунков и тд (вместо ф-л что-то нужно придумать)


\usepackage{subcaption}
\usepackage{circuitikz}

% Числа на осях: десятичная запятая, без разделителя тысяч
\pgfplotsset{/pgf/number format/.cd, use comma, 1000 sep={}}

% --- Служебные макросы ------------------------------------------------------
\newcommand{\qq}{\textbf{???}}
\newcommand{\fillme}{\textbf{?? НЕОБХОДИМО ЗАПОЛНИТЬ ??}}
\newcommand{\filename}[1]{\texttt{\detokenize{#1}}}
% Итоговые числа из data/processed/results.tex; нет значения -> ???
\IfFileExists{../data/processed/results.tex}{% Автоматически создано scripts/process_data.py -- вручную не править.
\expandafter\def\csname res@Mx\endcsname{29{,}8 \pm 0{,}5}
\expandafter\def\csname res@epsMx\endcsname{1{,}8}
\expandafter\def\csname res@MxNom\endcsname{33{,}3}
\expandafter\def\csname res@MxDev\endcsname{11}
\expandafter\def\csname res@R0eff\endcsname{0{,}336 \pm 0{,}006}
\expandafter\def\csname res@R0\endcsname{0{,}30}
\expandafter\def\csname res@Kyp\endcsname{38{,}4 \pm 0{,}7}
\expandafter\def\csname res@KyNom\endcsname{40{,}0}
\expandafter\def\csname res@KyDev\endcsname{4{,}0}
\expandafter\def\csname res@epsKy\endcsname{1{,}8}
\expandafter\def\csname res@UinPP\endcsname{6{,}00}
\expandafter\def\csname res@UoutPP\endcsname{54{,}0}
\expandafter\def\csname res@ratioU\endcsname{111}
\expandafter\def\csname res@tauExp\endcsname{0{,}354 \pm 0{,}017}
\expandafter\def\csname res@tauRC\endcsname{0{,}40}
\expandafter\def\csname res@tauDev\endcsname{12}
\expandafter\def\csname res@omegaTau\endcsname{126}
\expandafter\def\csname res@omegaTauExp\endcsname{111}
\expandafter\def\csname res@Rint\endcsname{20}
\expandafter\def\csname res@Cint\endcsname{20}
\expandafter\def\csname res@h3\endcsname{4{,}46}
\expandafter\def\csname res@b3\endcsname{0{,}674}
\expandafter\def\csname res@Hs3\endcsname{71{,}4 \pm 2{,}6}
\expandafter\def\csname res@Bs3\endcsname{1{,}43 \pm 0{,}05}
\expandafter\def\csname res@Hc3\endcsname{29{,}0 \pm 2{,}3}
\expandafter\def\csname res@Br3\endcsname{1{,}39 \pm 0{,}05}
\expandafter\def\csname res@ratio3\endcsname{0{,}97}
\expandafter\def\csname res@scrHs3\endcsname{49{,}1 \pm 1{,}4}
\expandafter\def\csname res@scrBs3\endcsname{1{,}43 \pm 0{,}05}
\expandafter\def\csname res@scrHc3\endcsname{25{,}0 \pm 0{,}6}
\expandafter\def\csname res@scrBr3\endcsname{1{,}348 \pm 0{,}029}
\expandafter\def\csname res@scrratio3\endcsname{0{,}94}
\expandafter\def\csname res@muMax3\endcsname{(6{,}0 \pm 2{,}9) \cdot 10^{4}}
\expandafter\def\csname res@HmuMax3\endcsname{16}
\expandafter\def\csname res@Hend3\endcsname{232}
\expandafter\def\csname res@Bend3\endcsname{1{,}43}
}{}
\newcommand{\res}[1]{\ifcsname res@#1\endcsname\csname res@#1\endcsname\else\qq\fi}
% Фото/зарисовка из report/img/. Если файла ещё нет -- рамка-заглушка.
\newcommand{\labimage}[2]{%
    \IfFileExists{img/#1}{%
        \includegraphics[width=#2\linewidth]{img/#1}%
    }{%
        \fbox{\parbox[c][0.6\dimexpr#2\linewidth\relax][c]{%
                \dimexpr#2\linewidth-2\fboxsep-2\fboxrule\relax}{%
            \centering\small\filename{img/#1}\\[2mm]\fillme}}%
    }%
}
\newcommand{\vx}{\text{вх}}
\newcommand{\vy}{\text{вых}}

% Колонтитулы (как в остальных работах репозитория)
\pagestyle{fancy}
\fancyhf{}
\fancyhead[L]{%
    \IfFileExists{../../common/latex_template/miptlogo.png}{%
        \includegraphics[height=1cm]{../../common/latex_template/miptlogo.png}%
    }{%
        \textbf{[ЛОГОТИП МФТИ]}%
    }%
}
\fancyhead[R]{\textbf{Емельянов ДА, Байдин ФМ}}
\fancyfoot[C]{\thepage}
\renewcommand{\headrulewidth}{0.4pt}
\setlength{\headheight}{33pt}

\begin{document}

% --- ТИТУЛЬНЫЙ ЛИСТ ---
\begin{titlepage}
\begin{center}
    {\large МОСКОВСКИЙ ФИЗИКО-ТЕХНИЧЕСКИЙ ИНСТИТУТ}
\end{center}
\begin{center}
    {\large Физтех-школа физики и исследований им. Ландау\break ЛФИ}
\end{center}

    \vspace{3cm}

\vspace{0.15cm}
{
\begin{center}
    \huge \bf Лабораторная работа №3.4.5 по курсу Общая физика\\
    \vspace{0.4cm}
    \LARGE Петля гистерезиса (динамический метод).
\end{center}
}
\vspace{1.5cm}
\begin{flushright}
Емельянов Дмитрий Алексеевич\\
Байдин Фаддей Михайлович\\
\vspace{0.4cm}
Б02-502
\end{flushright}
\vspace{10cm}
\begin{center}
    \large \qq\ октября 2026, г. Долгопрудный
\end{center}
\end{titlepage}
% ----------------------

\section{Аннотация}
В работе с помощью электронного осциллографа в режиме $X$--$Y$ исследуются
петли гистерезиса нескольких тороидальных ферромагнитных образцов в
переменном поле частоты $\nu_0 = 50$~Гц. Сигнал, пропорциональный
напряжённости поля $H$, снимается с резистора в цепи намагничивающей
обмотки, а сигнал, пропорциональный индукции $B$, --- с интегрирующей
$RC$-цепочки, подключённой к измерительной обмотке.

Для каждого образца по предельной петле определяются коэрцитивное поле
$H_c$, остаточная индукция $B_r$, индукция насыщения $B_s$ и
соответствующая ей амплитуда поля $H_s$; по вершинам частных петель
строится начальная кривая намагничивания и оцениваются начальная и
максимальная дифференциальная магнитные проницаемости. Проводятся
калибровка каналов осциллографа и измерение постоянной времени
интегрирующей цепочки.

\section{Теоретические сведения}

\subsection{Петля гистерезиса ферромагнетика}
Индукция $B$ и напряжённость $H$ магнитного поля в ферромагнетике связаны
неоднозначно: $B$ зависит не только от $H$, но и от предыстории образца
(\cref{fig:loop_theory}). При монотонном увеличении поля из размагниченного
состояния ($H = 0$, $B = 0$) образец намагничивается по
\emph{начальной кривой намагничивания} $OAC$; в точке $C$ намагниченность
достигает насыщения $M_s$, и далее (участок $CD$) индукция растёт только за
счёт роста $H$:
\begin{equation}
    \label{eq:B_sat}
    B = \mu_0 (H + M_s), \qquad \frac{dB}{dH} = \mu_0 .
\end{equation}

\begin{figure}[!ht]
    \centering
    \begin{tikzpicture}
        \begin{axis}[
            width=0.62\textwidth, height=0.48\textwidth,
            axis lines=middle, xlabel={$H$}, ylabel={$B$},
            xmin=-4.2, xmax=4.6, ymin=-1.55, ymax=1.65,
            xtick=\empty, ytick=\empty, samples=200, clip=false,
        ]
            % предельная петля: B = tanh((H -+ Hc)/a) + малый линейный рост
            \addplot[thick, domain=-3.6:3.6] {tanh((x+0.9)/0.55) + 0.05*x};
            \addplot[thick, domain=-3.6:3.6] {tanh((x-0.9)/0.55) + 0.05*x};
            % начальная кривая
            \addplot[thick, dashed, domain=0:3.6] {tanh(x*x/1.6) + 0.05*x};
            % частная петля (пунктир)
            \addplot[densely dotted, domain=-1.3:1.3] {0.86*tanh((x+0.35)/0.45)};
            \addplot[densely dotted, domain=-1.3:1.3] {0.86*tanh((x-0.35)/0.45)};
            \node[above left]  at (axis cs:3.6,1.18) {$D$};
            \node[above left]  at (axis cs:2.4,1.1) {$C$};
            \node[below right] at (axis cs:-2.4,-1.1) {$C'$};
            \node[below right] at (axis cs:-3.6,-1.18) {$D'$};
            \node[below right] at (axis cs:0,0.95) {$B_r$};
            \fill (axis cs:0,0.934) circle (1.2pt) (axis cs:-0.9,0) circle (1.2pt) (axis cs:1.3,0.85) circle (1.2pt) (axis cs:-1.3,-0.85) circle (1.2pt);
            \node[below left]  at (axis cs:-0.9,0) {$-H_c$};
            \node[right] at (axis cs:1.35,0.78) {$A$};
            \node[left] at (axis cs:-1.35,-0.78) {$A'$};
            \node[below right] at (axis cs:0,0) {$O$};
            \draw[gray, dashed] (axis cs:0,1.0) -- (axis cs:3.6,1.18);
            \node[left] at (axis cs:0,1.08) {\small $\mu_0 M_s$};
        \end{axis}
    \end{tikzpicture}
    \caption{Петля гистерезиса ферромагнетика (схематично): $CEFC'E'F'C$ ---
             предельная петля, $OAC$ --- начальная кривая намагничивания,
             $AA'$ --- частная петля}
    \label{fig:loop_theory}
\end{figure}

Если к образцу прикладывать переменное поле, его состояние на плоскости
$H$--$B$ описывает замкнутую кривую --- \emph{петлю гистерезиса}. Если
амплитуда поля достаточна для насыщения, получается \emph{предельная петля};
при дальнейшем увеличении амплитуды растут только «усы» $CD$ и $C'D'$, а
площадь петли практически не меняется. Пересечение предельной петли с осью
$B$ даёт \emph{остаточную индукцию} $B_r$, с осью $H$ --- \emph{коэрцитивное
поле} $H_c$. Вершины частных петель (точки $A$, $A'$) лежат на начальной
кривой намагничивания, поэтому её можно снять, плавно уменьшая амплитуду
тока и отмечая положения вершин.

Наклон кривой намагничивания характеризуется дифференциальной магнитной
проницаемостью
\begin{equation}
    \label{eq:mu_diff}
    \mu_{\text{диф}} = \frac{1}{\mu_0}\frac{dB}{dH};
\end{equation}
начальная проницаемость $\mu_{\text{нач}} = \lim\limits_{H\to 0} B/(\mu_0 H)$.

Площадь петли на плоскости $H$--$B$ равна энергии, которая необратимо
выделяется в виде тепла в единице объёма образца за один цикл
перемагничивания:
\begin{equation}
    \label{eq:loss}
    w = \oint H\,dB, \qquad \frac{P}{V} = \nu_0\, w .
\end{equation}

\subsection{Измерение напряжённости поля}
Для тороидального образца с намагничивающей обмоткой из $N_0$ витков по
теореме о циркуляции
\begin{equation}
    \label{eq:H}
    H = \frac{N_0 I_0}{2\pi R} = \frac{N_0 I_0}{L},
\end{equation}
где $L = 2\pi R$ --- длина средней линии тора (намагниченность можно
считать однородной при $r \ll R$). Ток $I_0$ измеряется по падению
напряжения $U_R = R_0 I_0$ на резисторе $R_0$, которое подаётся на канал $X$
осциллографа. Поэтому отклонение луча по горизонтали пропорционально $H$:
\begin{equation}
    \label{eq:hx}
    H = h_x\,[X], \qquad h_x = \frac{N_0}{L}\, M_x, \qquad
    M_x = \frac{K_x}{R_0},
\end{equation}
где $K_x$ --- коэффициент отклонения канала $X$ (мВ/дел), $M_x$ --- цена
деления по току (мА/дел).

Цену деления $M_x$ удобно определить калибровкой: при закороченной обмотке
$N_0$ ток синусоидален, и по показанию амперметра $I_{\text{эфф}}$ полная
длина горизонтальной развёртки $[2X]$ соответствует удвоенной амплитуде
тока:
\begin{equation}
    \label{eq:calib_x}
    M_x = \frac{2\sqrt2\, I_{\text{эфф}}}{[2X]} .
\end{equation}
(С образцом калибровать нельзя: катушка с ферромагнетиком --- нелинейный
элемент, и ток в ней несинусоидален.)

\subsection{Измерение индукции. Интегрирующая цепочка}
Если измерительная обмотка из $N_{\text{и}}$ витков плотно охватывает образец
сечения $S$, то по закону электромагнитной индукции
\begin{equation}
    \label{eq:B_emf}
    |B| = \frac{1}{SN_{\text{и}}}\int \mathcal{E}\,dt .
\end{equation}
Для интегрирования используется $RC$-цепочка (\cref{fig:rc}): ЭДС подаётся
на последовательно соединённые $R_{\text{и}}$ и $C_{\text{и}}$, а
напряжение снимается с конденсатора. Если сопротивление источника мало,
входное сопротивление осциллографа велико, а $U_\vy \ll U_\vx$, то
$I \approx U_\vx / R_{\text{и}}$ и
\begin{equation}
    \label{eq:integrator}
    U_\vy = \frac{q}{C_{\text{и}}}
    \approx \frac{1}{\tau_{\text{и}}}\int_0^t U_\vx\,dt,
    \qquad \tau_{\text{и}} = R_{\text{и}} C_{\text{и}} .
\end{equation}
Отсюда
\begin{equation}
    \label{eq:B}
    |B| = \frac{\tau_{\text{и}}}{SN_{\text{и}}}\, U_C, \qquad
    B = b_y\,[Y], \qquad b_y = \frac{\tau_{\text{и}} K_y}{SN_{\text{и}}},
\end{equation}
где $K_y$ --- коэффициент отклонения канала $Y$.

\begin{figure}[!ht]
    \centering
    \begin{circuitikz}
        \draw (0,0) to[short, o-] (0.5,0) to[R, l=$R_{\text{и}}$] (3,0)
              to[short, -o] (4.5,0);
        \draw (3,0) to[C, l=$C_{\text{и}}$] (3,-2);
        \draw (0,-2) to[short, o-o] (4.5,-2);
        \draw (0,0) to[open, v=$U_\vx$] (0,-2);
        \draw (4.5,0) to[open, v^=$U_\vy$] (4.5,-2);
    \end{circuitikz}
    \caption{Интегрирующая $RC$-ячейка}
    \label{fig:rc}
\end{figure}

Для синусоидального сигнала частоты $\omega_0$ отношение амплитуд
\begin{equation}
    \frac{U_\vy}{U_\vx} = \frac{1/(\omega_0 C)}{\sqrt{R^2 + 1/(\omega_0C)^2}},
\end{equation}
так что условие применимости \cref{eq:integrator}
\begin{equation}
    \label{eq:tau_cond}
    \omega_0\tau_{\text{и}} \gg 1,
\end{equation}
и тогда
\begin{equation}
    \label{eq:tau}
    \frac{U_\vy}{U_\vx} \approx \frac{1}{\omega_0\tau_{\text{и}}}
    \quad\Longrightarrow\quad
    \tau_{\text{и}} = \frac{U_\vx}{\omega_0 U_\vy} .
\end{equation}
Если $\omega_0$ --- самая низкая гармоника несинусоидального сигнала, то
при $\omega_0\tau \gg 1$ интегрирование корректно и для всех высших гармоник.

\section{Используемое оборудование}
\begin{itemize}
    \item трансформаторный блок: регулировочный автотрансформатор и
          разделительный понижающий трансформатор (сеть 220~В, 50~Гц);
    \item тороидальные образцы с намагничивающей ($N_0$) и измерительной
          ($N_{\text{и}}$) обмотками: \fillme;
    \item резистор $R_0 = \res{R0}$~Ом в цепи намагничивающей обмотки;
    \item интегрирующая $RC$-цепочка: $R_{\text{и}} = \res{Rint}$~кОм,
          $C_{\text{и}} = \res{Cint}$~мкФ;
    \item амперметр $A$ (действующее значение тока): \fillme;
    \item вольтметр $V$, делитель напряжения 1:10, 1:100;
    \item электронный осциллограф: \fillme.
\end{itemize}

\section{Экспериментальная установка и методика измерений}
Схема установки показана на \cref{fig:setup}. Напряжение сети через
трансформаторный блок $T$ подаётся на намагничивающую обмотку $N_0$
образца, последовательно с которой включены амперметр $A$ и резистор $R_0$.
Напряжение $U_R = R_0 I_0 \propto H$ подаётся на канал $X$ осциллографа.
ЭДС измерительной обмотки $N_{\text{и}}$ интегрируется $RC$-цепочкой, и
напряжение $U_C \propto B$ подаётся на канал $Y$. В режиме $X$--$Y$ на
экране наблюдается петля гистерезиса в масштабах $h_x$, $b_y$
(\cref{eq:hx,eq:B}).

\begin{figure}[!ht]
    \centering
    \begin{circuitikz}[scale=0.95, transform shape]
        % первичная цепь
        \draw (0,0) to[sV, l_=$U_0$] (0,3)
              to[rmeter, t=A] (2.6,3)
              to[short] (3.4,3)
              to[cute inductor, l_=$N_0$] (3.4,0)
              to[R, l_=$R_0$] (0,0);
        \node[left, align=center] at (-1.1,1.5) {\small $T$\\[-1mm]\small $\approx$220~В,\\[-1mm]\small 50~Гц};
        \draw (3.4,0) to[short, *-o] (3.4,-1.1) node[below] {\small $X$};
        \draw (0,0) to[short, *-o] (0,-1.1) node[below] {\small общий};
        % сердечник
        \draw[very thick] (4.05,0.2) -- (4.05,2.8);
        \draw[very thick] (4.25,0.2) -- (4.25,2.8);
        \node at (4.15,3.35) {\small образец};
        % измерительная цепь
        \draw (4.9,3) to[cute inductor, l=$N_{\text{и}}$] (4.9,0);
        \draw (4.9,3) to[R, l=$R_{\text{и}}$] (8,3) to[C, l=$C_{\text{и}}$] (8,0) -- (4.9,0);
        \draw (8,3) to[short, *-o] (9.4,3) node[right] {\small $Y$};
        \draw (8,0) to[short, *-o] (9.4,0) node[right] {\small общий};
    \end{circuitikz}
    \caption{Схема установки (упрощённо; цепь калибровки канала $Y$ через
             делитель и вольтметр $V$ не показана)}
    \label{fig:setup}
\end{figure}

Порядок измерений:
\begin{enumerate}
    \item Подбирается ток в намагничивающей обмотке, при котором наблюдается
          предельная петля (при изменении тока меняется лишь длина «усов»).
          Коэффициенты $K_x$, $K_y$ выбираются так, чтобы петля занимала
          б\'{о}льшую часть экрана (ручки плавной регулировки --- в
          положении калибровки). Петля фотографируется.
    \item По сетке экрана измеряются полные ширина и высота петли $[2X_s]$,
          $[2Y_s]$, ширина петли на оси абсцисс $[2X_c]$ и высота на оси
          ординат $[2Y_r]$.
    \item Плавно уменьшая ток до нуля, отмечаются положения вершин
          $8\div10$ частных петель $[2X_A]$, $[2Y_A]$ --- точки начальной
          кривой намагничивания.
    \item Измерения повторяются для всех образцов.
    \item Калибровка канала $X$ по \cref{eq:calib_x} (обмотка $N_0$
          закорочена) для всех использованных $K_x$; проверка канала $Y$
          подачей известного напряжения через делитель 1:100.
    \item Измерение отношения $U_\vx/U_\vy$ для $RC$-ячейки и расчёт
          $\tau_{\text{и}}$ по \cref{eq:tau}.
\end{enumerate}
Погрешность отсчёта по сетке экрана принята равной $\pm0{,}1$~дел.

\section{Результаты измерений и обработка данных}

\subsection{Параметры образцов}
Параметры тороидов (с установки) приведены в \cref{tab:samples}.

\begin{table}[!ht]
    \centering
    \caption{Параметры образцов}
    \label{tab:samples}
    \csvreader[
        tabularray = {
            colspec = {|c|c|c|c|c|c|},
            hlines, vlines, row{1} = {font=\bfseries},
        },
        table head = {
            № & Материал & $N_0$ & $N_{\text{и}}$ & $S$, см$^2$ & $2\pi R$, см \\
        },
    ]%
    {../data/processed/samples_processed.csv}%
    {sample=\cS, material=\cM, N0=\cN, N_i=\cNi, S_cm2=\cSs, L_cm=\cL}%
    {\cS & \cM & $\cN$ & $\cNi$ & $\cSs$ & $\cL$}
\end{table}

\subsection{Калибровка канала $X$}
Результаты калибровки при закороченной обмотке $N_0$ приведены в
\cref{tab:calib_x}; $M_x^{\text{эксп}}$ рассчитано по \cref{eq:calib_x},
$M_x^{\text{теор}} = K_x/R_0$.

\begin{table}[!ht]
    \centering
    \caption{Калибровка канала $X$}
    \label{tab:calib_x}
    \csvreader[
        tabularray = {
            colspec = {|c|c|c|c|c|c|},
            hlines, vlines, row{1} = {font=\bfseries},
        },
        table head = {
            $K_x$, мВ/дел & $I_{\text{эфф}}$, мА & $[2X]$, дел &
            $M_x^{\text{эксп}}$, мА/дел & $M_x^{\text{теор}}$, мА/дел &
            $\varepsilon$, \% \\
        },
    ]%
    {../data/processed/calib_x_processed.csv}%
    {Kx=\cK, I_eff_mA=\cI, X2=\cX, Mx_exp=\cMe, Mx_th=\cMt, dev_percent=\cD}%
    {$\cK$ & $\cI$ & $\cX$ & $\cMe$ & $\cMt$ & $\cD$}
\end{table}

Эффективное сопротивление, определённое по калибровке,
$R_0^{\text{эфф}} = K_x/M_x = \res{R0eff}$~Ом (номинал $R_0 = \res{R0}$~Ом).
Относительная погрешность цены деления по току
$\varepsilon_{M_x} \approx \res{epsMx}$~\% (класс амперметра и отсчёт длины
развёртки).

\subsection{Проверка калибровки канала $Y$}
\begin{table}[!ht]
    \centering
    \caption{Проверка канала $Y$: $M_y = 2\sqrt2\,U_{\text{эфф}}/[2Y]$}
    \label{tab:calib_y}
    \csvreader[
        tabularray = {
            colspec = {|c|c|c|c|c|},
            hlines, vlines, row{1} = {font=\bfseries},
        },
        table head = {
            $K_y$, мВ/дел & $U_{\text{эфф}}$, мВ & $[2Y]$, дел &
            $M_y$, мВ/дел & $\varepsilon$, \% \\
        },
    ]%
    {../data/processed/calib_y_processed.csv}%
    {Ky=\cK, U_eff_mV=\cU, Y2=\cY, My_exp=\cM, dev_percent=\cD}%
    {$\cK$ & $\cU$ & $\cY$ & $\cM$ & $\cD$}
\end{table}

Измеренная цена деления отличается от номинальной $K_y$ не более чем на
$\res{devY}$~\% (\cref{tab:calib_y}), что согласуется с паспортной
погрешностью осциллографа; в расчётах используются номинальные значения
$K_y$.

\subsection{Постоянная времени интегрирующей цепочки}
\begin{table}[!ht]
    \centering
    \caption{Измерение постоянной времени $RC$-цепочки}
    \label{tab:rc}
    \csvreader[
        tabularray = {
            colspec = {|c|c|c|c|},
            hlines, vlines, row{1} = {font=\bfseries},
        },
        table head = {
            $U_\vx$, мВ & $U_\vy$, мВ & $\nu$, Гц & $\tau_{\text{и}}$, мс \\
        },
    ]%
    {../data/processed/rc_processed.csv}%
    {U_in_mV=\cA, U_out_mV=\cB, nu_Hz=\cN, tau_ms=\cT}%
    {$\cA$ & $\cB$ & $\cN$ & $\cT$}
\end{table}

По \cref{eq:tau}
\begin{equation*}
    \tau_{\text{и}}^{\text{эксп}} = (\res{tauExp})~\text{мс},
    \qquad
    \tau_{\text{и}}^{RC} = R_{\text{и}}C_{\text{и}} = \res{tauRC}~\text{мс};
\end{equation*}
расхождение составляет $\res{tauDev}$~\%. Условие \cref{eq:tau_cond}
выполняется: $\omega_0\tau_{\text{и}} \approx \res{omegaTau} \gg 1$.
В дальнейших расчётах используется $\tau_{\text{и}}^{\text{эксп}}$.

\subsection{Предельные петли гистерезиса}
Фотографии предельных петель приведены на \cref{fig:loops}.

\begin{figure}[!ht]
    \centering
    \begin{subfigure}{0.32\textwidth}
        \centering \labimage{loop_1.jpg}{1.0}
        \caption{образец 1 (\res{mat1})}
    \end{subfigure}
    \hfill
    \begin{subfigure}{0.32\textwidth}
        \centering \labimage{loop_2.jpg}{1.0}
        \caption{образец 2 (\res{mat2})}
    \end{subfigure}
    \hfill
    \begin{subfigure}{0.32\textwidth}
        \centering \labimage{loop_3.jpg}{1.0}
        \caption{образец 3 (\res{mat3})}
    \end{subfigure}
    \caption{Предельные петли гистерезиса на экране осциллографа}
    \label{fig:loops}
\end{figure}

Пересчёт делений в поля выполнен по \cref{eq:hx,eq:B}; для каждого образца
$H_s = h_x[2X_s]/2$, $B_s = b_y[2Y_s]/2$, $H_c = h_x[2X_c]/2$,
$B_r = b_y[2Y_r]/2$. Результаты сведены в \cref{tab:loops}. Погрешности
включают погрешность отсчёта по экрану, калибровки канала $X$,
коэффициента $K_y$ и постоянной времени $\tau_{\text{и}}$:
\begin{equation*}
    \varepsilon_H = \sqrt{\varepsilon_{M_x}^2 + \left(\frac{\Delta[X]}{[2X]}\right)^2},
    \qquad
    \varepsilon_B = \sqrt{\varepsilon_{K_y}^2 + \varepsilon_{\tau}^2
                          + \left(\frac{\Delta[Y]}{[2Y]}\right)^2}.
\end{equation*}

\begin{table}[!ht]
    \centering
    \small
    \caption{Характеристики предельных петель: $h_x$ и $b_y$ --- цены
             деления при измерении $[2X_s]$, $[2Y_s]$}
    \label{tab:loops}
    \csvreader[
        tabularray = {
            colspec = {|c|c|c|c|c|c|c|c|},
            hlines, vlines, row{1} = {font=\bfseries},
        },
        table head = {
            № & Материал & $h_x$, $\frac{\text{А/м}}{\text{дел}}$ &
            $b_y$, $\frac{\text{Тл}}{\text{дел}}$ &
            $H_s$, А/м & $B_s$, Тл & $H_c$, А/м & $B_r$, Тл \\
        },
    ]%
    {../data/processed/loops_processed.csv}%
    {sample=\cS, material=\cM, hx=\cH, by=\cB, Hs=\cHs, dHs=\cdHs,
     Bs=\cBs, dBs=\cdBs, Hc=\cHc, dHc=\cdHc, Br=\cBr, dBr=\cdBr}%
    {\cS & \cM & $\cH$ & $\cB$ & $\cHs\pm\cdHs$ & $\cBs\pm\cdBs$
     & $\cHc\pm\cdHc$ & $\cBr\pm\cdBr$}
\end{table}

Потери на перемагничивание оценены по площади петли $A$ (в дел$^2$,
по фотографии): $w = A\,h_x b_y$, $P/V = \nu_0 w$ (\cref{eq:loss}):
\begin{center}
\begin{tblr}{colspec={|c|c|c|c|}, hlines, vlines, row{1}={font=\bfseries}}
    Образец & 1 & 2 & 3 \\
    $w$, Дж/м$^3$ & $\res{w1}$ & $\res{w2}$ & $\res{w3}$ \\
    $P/V$, кВт/м$^3$ & $\res{P1}$ & $\res{P2}$ & $\res{P3}$ \\
\end{tblr}
\end{center}

\clearpage
\subsection{Начальные кривые намагничивания}
Вершины частных петель, пересчитанные в $H$ и $B$, показаны на
\cref{fig:initial}; на \cref{fig:mudiff} --- дифференциальная
проницаемость $\mu_{\text{диф}} = \Delta B/(\mu_0\Delta H)$, рассчитанная по
соседним точкам.

\begin{figure}[!ht]
    \centering
    \begin{tikzpicture}
        \begin{axis}[
            width=0.85\textwidth, height=0.55\textwidth,
            xlabel={$H$, А/м}, ylabel={$B$, Тл},
            xmin=0, ymin=0,
            grid=both, minor tick num=1,
            legend pos=south east,
            scaled x ticks=false,
        ]
            \addplot[thick, mark=*] table[col sep=comma, x=H, y=B]
                {../data/processed/initial_1.csv};
            \addlegendentry{образец 1}
            \addplot[thick, mark=square*] table[col sep=comma, x=H, y=B]
                {../data/processed/initial_2.csv};
            \addlegendentry{образец 2}
            \addplot[thick, mark=triangle*] table[col sep=comma, x=H, y=B]
                {../data/processed/initial_3.csv};
            \addlegendentry{образец 3}
        \end{axis}
    \end{tikzpicture}
    \caption{Начальные кривые намагничивания $B(H)$ (по вершинам частных петель)}
    \label{fig:initial}
\end{figure}

\begin{figure}[!ht]
    \centering
    \begin{tikzpicture}
        \begin{axis}[
            width=0.85\textwidth, height=0.5\textwidth,
            xlabel={$H$, А/м}, ylabel={$\mu_{\text{диф}}$},
            xmin=0, ymin=0,
            grid=both, minor tick num=1,
            legend pos=north east,
            scaled x ticks=false,
        ]
            \addplot[thick, mark=*] table[col sep=comma, x=H, y=mu_diff]
                {../data/processed/mudiff_1.csv};
            \addlegendentry{образец 1}
            \addplot[thick, mark=square*] table[col sep=comma, x=H, y=mu_diff]
                {../data/processed/mudiff_2.csv};
            \addlegendentry{образец 2}
            \addplot[thick, mark=triangle*] table[col sep=comma, x=H, y=mu_diff]
                {../data/processed/mudiff_3.csv};
            \addlegendentry{образец 3}
        \end{axis}
    \end{tikzpicture}
    \caption{Дифференциальная магнитная проницаемость вдоль начальной кривой}
    \label{fig:mudiff}
\end{figure}

\begin{table}[!ht]
    \centering
    \caption{Начальная и максимальная дифференциальная проницаемости}
    \label{tab:mu}
    \begin{tblr}{colspec={|c|c|c|c|}, hlines, vlines, row{1}={font=\bfseries}}
        Образец & $\mu_{\text{нач}}$ & $\mu_{\text{диф}}^{\max}$ &
            $H$ при $\mu_{\text{диф}}^{\max}$, А/м \\
        1 & $\res{muInit1}$ & $\res{muMax1}$ & $\res{HmuMax1}$ \\
        2 & $\res{muInit2}$ & $\res{muMax2}$ & $\res{HmuMax2}$ \\
        3 & $\res{muInit3}$ & $\res{muMax3}$ & $\res{HmuMax3}$ \\
    \end{tblr}
\end{table}

Начальная проницаемость оценена как $B/(\mu_0H)$ в первой (наименьшей по
$H$) точке; поскольку частные петли при малых токах плохо различимы на
экране, эта оценка грубая (погрешность порядка $20\div30$~\%).
\fillme: описать вид кривых (участок быстрого роста, выход на насыщение).

\section{Обсуждение результатов}
Сравнение с типичными (справочными) значениями приведено в
\cref{tab:compare}. Справочные значения относятся, как правило, к
квазистатическому перемагничиванию; при частоте 50~Гц из-за вихревых токов
петля «расширяется», поэтому динамическое $H_c$ обычно больше
статического.

\begin{table}[!ht]
    \centering
    \caption{Сравнение со справочными данными (порядок величины)}
    \label{tab:compare}
    \begin{tblr}{colspec={|l|c|c|c|c|}, hlines, vlines, row{1}={font=\bfseries}}
        Материал & $H_c^{\text{эксп}}$, А/м & $H_c^{\text{спр}}$, А/м
                 & $B_s^{\text{эксп}}$, Тл & $B_s^{\text{спр}}$, Тл \\
        \res{mat1} & $\res{Hc1}$ & \qq & $\res{Bs1}$ & \qq \\
        \res{mat2} & $\res{Hc2}$ & \qq & $\res{Bs2}$ & \qq \\
        \res{mat3} & $\res{Hc3}$ & \qq & $\res{Bs3}$ & \qq \\
    \end{tblr}
\end{table}

Ориентиры: феррит марганец-цинковый (2000НМ) --- $H_c \sim 10\div25$~А/м,
$B_s \approx 0{,}35\div0{,}4$~Тл; пермаллой (79НМ) --- $H_c \sim 1\div5$~А/м,
$B_s \approx 0{,}7\div0{,}8$~Тл; кремнистое (трансформаторное) железо ---
$H_c \sim 10\div50$~А/м, $B_s \approx 1{,}5\div2$~Тл.

\begin{itemize}
    \item \fillme: какой из образцов магнитомягкий (наименьшее $H_c$),
          какой обладает наибольшей индукцией насыщения и наибольшими
          потерями.
    \item Калибровка канала $X$ даёт $R_0^{\text{эфф}} = \res{R0eff}$~Ом,
          что \fillme\ с номиналом $R_0 = \res{R0}$~Ом.
    \item Постоянная времени, измеренная по \cref{eq:tau}, совпадает с
          $R_{\text{и}}C_{\text{и}}$ с точностью $\res{tauDev}$~\%;
          условие интегрирования $\omega_0\tau \gg 1$ выполнено с запасом
          $\omega_0\tau \approx \res{omegaTau}$.
    \item Основные источники погрешности: отсчёт по сетке экрана
          (особенно для узких петель и малых частных петель), класс
          амперметра, точность коэффициентов отклонения осциллографа.
          Систематически: $U_\vy$ не пренебрежимо мало по сравнению с
          $U_\vx$ (относительная поправка $\sim 1/(\omega_0\tau)^2$),
          неоднородность поля по сечению тора ($H \propto 1/r$),
          скин-эффект в проводящих образцах.
\end{itemize}

\section{Заключение}
В работе динамическим методом (на частоте 50~Гц) исследованы петли
гистерезиса трёх ферромагнитных образцов:
\begin{itemize}
    \item образец 1 (\res{mat1}): $H_c = (\res{Hc1})$~А/м,
          $B_r = (\res{Br1})$~Тл, $B_s = (\res{Bs1})$~Тл при
          $H_s = (\res{Hs1})$~А/м;
    \item образец 2 (\res{mat2}): $H_c = (\res{Hc2})$~А/м,
          $B_r = (\res{Br2})$~Тл, $B_s = (\res{Bs2})$~Тл при
          $H_s = (\res{Hs2})$~А/м;
    \item образец 3 (\res{mat3}): $H_c = (\res{Hc3})$~А/м,
          $B_r = (\res{Br3})$~Тл, $B_s = (\res{Bs3})$~Тл при
          $H_s = (\res{Hs3})$~А/м;
    \item по начальным кривым намагничивания получены максимальные
          дифференциальные проницаемости $\mu_{\text{диф}}^{\max} \approx
          \res{muMax1}$, $\res{muMax2}$, $\res{muMax3}$ соответственно;
    \item постоянная времени интегрирующей цепочки
          $\tau_{\text{и}} = (\res{tauExp})$~мс, условие
          $\omega_0\tau \gg 1$ выполнено.
\end{itemize}
\fillme: вывод о соответствии справочным данным.

\end{document}
